


Generated by \texttt{python -m crashguard.vf2\_split\_doc} from \texttt{results/results.json} (keys \texttt{vz.v4.*}, \texttt{vf2.split.*}). Labels: MEASURED\_SOFTWARE, DESCRIPTIVE (no hypothesis is tested here). Delta verification item D7.4.

\subsection*{1. What 0.901 is}

The value 0.901 is the system-level recall of the frozen 15-tree forest on the VZCrash \textbf{test split as supplied} (mean over seeds 0 to 4; 95\% interval 0.892 to 0.910): the share of all crash events in the supplied test split whose score reaches the validation threshold at the maximum of F1. The unit of the split in the released files is the event (one clip, one trigger-aligned window); the dataset authors state that the split separates vehicles, but the release has no vehicle or trip identifier, so that statement cannot be checked here. The supplied test split contains events whose raw accelerometer array is identical to a train or validation event, and events whose trigger-aligned window equals a train or validation window. \textbf{It is not the recall of a naive window-level split}, and it is not a vehicle-grouped result. A reader who takes 0.901 as the recall of a split that was cleaned of duplicates reads it wrongly: it is the recall with the duplicates in.

\subsection*{2. The audit as a reusable procedure}

Input: a dataset with a supplied split into train, validation and test, one array of sensor samples per event. Output: the cross-split duplicate sets and the effect of removing them on a fixed detector. Code: \texttt{src/crashguard/vz\_audit.py} (independent re-derivation from the raw files; it does not import the pipeline modules) and \texttt{src/crashguard/vz\_leakage.py} (the test-set versions).

\begin{itemize}
\item \textbf{Exact array hashing.} SHA-256 over the bytes of the full accelerometer array of each event (and, separately, over the full event: accelerometer, gyroscope and speed). Group events by hash. A group that spans more than one split is a cross-split duplicate set; every test event in such a group is an exact copy of a train or validation event.
\item \textbf{Window-identity matching.} For each triggered event cut the trigger-aligned window (1.0 s before and 1.8 s after the earliest sample whose magnitude exceeds 2 g, 280 samples at 100 Hz; every axis clipped to +-16 g beforehand; edge-repeat padding). A test window counts as an identical window if some train or validation window has a correlation above 0.999 (axis means removed, axes concatenated) \textbf{and} a maximum absolute difference below 0.01 g. Both thresholds must hold.
\item \textbf{The 16 g and trigger rules.} Events with a raw peak above 16 g are dropped, as in the pipeline (2 test events); events that never exceed the 2 g trigger are never windowed and are scored as negatives (-1), so they count as missed crashes at system level.
\item \textbf{Test-set versions.} Evaluate the same fitted detector, with the same thresholds chosen on validation, on (a) the test split as supplied, (b) without exact copies, (c) without exact copies and without identical windows. Intervals come from one bootstrap over the supplied test events, so the three versions are paired. The test split is scored once per model and seed.
\end{itemize}

Use on another dataset: replace the window rule (length, trigger, clip) by the dataset's own and keep the two thresholds as parameters. The thresholds 0.999 and 0.01 g were written down before the audit documents were produced (\texttt{docs/07\_log/DECISIONS.md}, 2026-11-12) and are not tuned.

\subsection*{3. Recall of the frozen 15-tree forest under four test-set versions}

{\scriptsize\raggedright
\begin{longtable}{p{0.136\linewidth}p{0.136\linewidth}p{0.136\linewidth}p{0.136\linewidth}p{0.136\linewidth}p{0.136\linewidth}p{0.136\linewidth}}
\toprule
Test-set version & Test events & Crash events & System recall at the F1 threshold (95\% interval) & Recall at 1\% false alarms (95\% interval) & Recall at 5\% false alarms & AP on triggered events \\
\midrule
\endhead
(a) as supplied & 24,172 & 4,221 & 0.901 (0.892 to 0.910) & 0.631 (0.618 to 0.645) & 0.935 & 0.934 \\
(b) exact copies removed & 24,054 & 4,107 & 0.900 (0.891 to 0.909) & 0.627 (0.613 to 0.641) & 0.934 & 0.931 \\
(c) exact copies and identical windows removed & 22,862 & 2,969 & 0.873 (0.861 to 0.884) & 0.558 (0.541 to 0.576) & 0.914 & 0.892 \\
(d) ILLUSTRATIVE: naive random re-split of the train and validation events (test split not read; 70/15/15 by event) & 24,767 & 4,049 & 0.897 (0.893 to 0.902) & 0.619 (0.612 to 0.626) & 0.929 & 0.927 \\
\bottomrule
\end{longtable}}

Rows (a) to (c) are the stored audit results (\texttt{vz.v4.rf15.*}). Removing exact copies changes the system recall by 0.0006 (interval -0.0008 to 0.0018); removing exact copies and identical windows changes it by 0.0280 (0.0235 to 0.0326) and the recall at 1\% false alarms by 0.0727 (0.0641 to 0.0809). Row (d) is built from the supplied train and validation events only (\texttt{src/crashguard/vf2\_split\_sens.py}, seeds 0 to 4, the same forest and operating points); it is illustrative, not a test of a hypothesis, and its parts hold different events from the supplied test split. The pooled events already contain the cross-split duplicate groups of the release, so a random re-split spreads them over the parts and the result lies close to row (a): the duplicates in the supplied test split inflate the recall about as much as a random re-split does. The cleaner comparison is between rows (a) and (c).

\subsection*{4. Why there is no vehicle-grouped split}

VZCrash has no vehicle or trip identifier (\texttt{docs/05\_data\_code/DATASET\_FACTS.md}; the only grouping field is the vehicle size class). A split that keeps one vehicle in one part therefore cannot be built, and no result here is a vehicle-grouped result. Where a split inside the training data was needed (the folds of L2), event-level folds were used and the near-duplicate risk was recorded (\texttt{docs/07\_log/DECISIONS.md}, 2026-11-15). The most independent grouping available is the group of identical arrays (the hash groups of step 1); a split that keeps each hash group in one part removes the exact copies by construction and is the closest substitute.

